\section{The nonzero Poisson frequencies}
\label{sec:nonzero}

We prove the estimate that determines the range of \Cref{thm:main}.  The
argument uses only finite completion and the classical Weil bound
\cite{Weil1948}; it does not correlate the two coefficient sequences.

For a dyadic block as in \Cref{sec:poisson}, the $t$-integral and the change
of variables in \eqref{eq:poisson-jacobian} must be kept together.  Along
$hn-klm=r$, the phase is $x/(x+r)$, and
\begin{equation}\label{eq:phase-derivative}
 x\frac{d}{dx}\left(t\log\frac{x}{x+r}\right)
 \ll\frac{T|r|}{S}\ll T^\eps.
\end{equation}
Consequently the Poisson integral has size $TS$, with normalized derivatives
bounded by $TS\,T^\eps$.  The factor $T$ here is essential.

For $d=(h,km)$ as in \eqref{eq:gcd-decomposition}, integration by parts in
the Poisson integral restricts a nonzero Fourier frequency to
\begin{equation}\label{eq:a-range}
 0<|a|\leq A_d,\qquad
 A_d\ll\frac{H_0K_0M_0}{dS}T^\eps
      \asymp\frac{S}{dL_0N_0}T^\eps.
\end{equation}
Also
\begin{equation}\label{eq:rho-range}
                  0<|\rho|\ll \frac{R_0}{d},\qquad
                  R_0=ST^{-1+\eps}.
\end{equation}
If either upper bound is less than one, the corresponding block is empty.

\begin{lemma}[Incomplete reciprocal sum]\label{lem:reciprocal-sum}
Let $q,c\geq1$, and let $F$ be supported in an interval of length $X$ with
$F^{(j)}\ll_j ZX^{-j}$.  Then
\begin{equation}\label{eq:reciprocal-sum}
 \sum_{(u,q)=1}F(u)e\!\left(\frac{c\overline u}{q}\right)
 \ll_\eps Zq^\eps
 \left\{
 Xq^{-1/2}(c,q)^{1/2}
 +q^{1/2}(c,q)^{1/2}
 \right\}.
\end{equation}
Moreover,
\begin{equation}\label{eq:gcd-average}
                  \sum_{q\asymp Q}(c,q)^{1/2}
                  \ll_\eps Qc^\eps.
\end{equation}
\end{lemma}

\begin{proof}
With the Fourier-transform convention compatible with $e(x)=e^{2\pi ix}$,
finite completion gives
\[
 \sum_{(u,q)=1}F(u)e\!\left(\frac{c\overline u}{q}\right)
 =\frac1q\sum_{n\in\mathbb Z}S(c,n;q)\widehat F(n/q),
\]
where
\[
 S(c,n;q)=\sum_{u\bmod q}^{*}
 e\!\left(\frac{c\overline u+nu}{q}\right).
\]
For every fixed $A\geq2$, repeated integration by parts gives
\[
 |\widehat F(\xi)|\ll_A ZX(1+X|\xi|)^{-A},
 \qquad
 \sum_{n\in\mathbb Z}|\widehat F(n/q)|\ll_A Z(X+q).
\]
The Weil bound
$|S(c,n;q)|\leq\tau(q)(c,n,q)^{1/2}q^{1/2}$, together with
$(c,n,q)^{1/2}\leq(c,q)^{1/2}$ for every $n$, therefore proves
\eqref{eq:reciprocal-sum}, uniformly for both $X\leq q$ and $X>q$.
Thus the factor $(c,q)^{1/2}$ occurs in both the $Xq^{-1/2}$ term and the
$q^{1/2}$ term.  The term $n=0$ is the Ramanujan sum and obeys the same
majorant.  Finally,
$(c,q)^{1/2}\leq\sum_{e\mid(c,q)}e^{1/2}$ gives
\[
 \sum_{q\asymp Q}(c,q)^{1/2}
 \leq\sum_{e\mid c}e^{1/2}\#\{q\asymp Q:e\mid q\}
 \ll Q\sum_{e\mid c}e^{-1/2}\ll_\eps Qc^\eps,
\]
which is \eqref{eq:gcd-average}.
\end{proof}

For the coefficient block put
\begin{equation}\label{eq:coefficient-block-norms}
 \mathcal A_{H_0}=\max_{h\asymp H_0}|a_h|,
 \qquad
 \mathcal B_{K_0}=\sum_{k\asymp K_0}|b_k|.
\end{equation}

\begin{proposition}[A fixed-orientation block estimate]
\label{prop:block-weil}
The contribution of the nonzero frequencies obtained by Poisson summation in
$l$ satisfies
\begin{equation}\label{eq:block-weil}
 E^{(l)}(H_0,K_0,L_0,M_0,N_0)
 \ll_\eps T^\eps\mathcal A_{H_0}\mathcal B_{K_0}
 H_0^{1/2}(H_0+M_0).
\end{equation}
\end{proposition}

\begin{proof}
Choose the dyadic partition using fixed functions
\[
 \omega_H,\omega_K,\omega_L,\omega_M,\omega_N
 \in C_c^\infty((1/2,3)).
\]
All their derivatives are bounded independently of the block parameters.
We first expose the smooth weight to which
\Cref{lem:reciprocal-sum} is applied.

Fix $d,a,\rho$ and $k$, and write
\begin{equation}\label{eq:d-k-split}
 g=(d,k),\qquad d=g\delta,\qquad k=gk_1,\qquad m=\delta m_1.
\end{equation}
The equality $(d,k)=g$ implies $(\delta,k_1)=1$: for every prime $p$,
the minimum of the $p$-adic valuations of $d/g$ and $k/g$ is zero.
Since $d\mid km$, it follows that $\delta\mid m$, so $m_1$ is an integer.
If $h=dq$, then
\begin{equation}\label{eq:exact-coprimality}
 (h,km)=d
 \quad\Longleftrightarrow\quad
 (q,k_1m_1)=1,
 \qquad \frac{km}{d}=k_1m_1.
\end{equation}
All substitutions are uniquely determined by $d,h,k,m$; conversely the
right side implies the original gcd condition.  Thus this is a bijective
parametrization, with neither omission nor repeated counting.

Put $r=d\rho$ and, after the change of variables $x=kml$ from
\eqref{eq:poisson-jacobian}, set
\begin{equation}\label{eq:completed-l-n-x}
 \ell_x=\frac{x}{km},\qquad n_x=\frac{x+d\rho}{h}.
\end{equation}
For the direct AFE piece define the exact normalized amplitude
\begin{align}
 \mathscr W^+_{h,k,m,d\rho}(x,t)
 ={}&
 \frac{(H_0K_0L_0M_0N_0)^{1/2}H_0K_0M_0}
 {hkm\,(hk\ell_xmn_x)^{1/2}}\notag\\
 &\times\omega_H(h/H_0)\omega_K(k/K_0)
 \omega_L(\ell_x/L_0)\notag\\
 &\times\omega_M(m/M_0)\omega_N(n_x/N_0)
 \ell_x^{-\alpha}m^{-\beta}n_x^{-\gamma}
 w(t)\left(\frac{x+d\rho}{x}\right)^{it}
 \notag\\
 &\times
 \V^+\!\left(\frac{x}{kmY},\frac{m(x+d\rho)}h,t\right).
 \label{eq:completed-normalized-amplitude}
\end{align}
The cutoffs make the integrand zero unless $x>0$ and $x+d\rho>0$.
The dual amplitude is obtained by the exact substitution
\eqref{eq:dual-substitution}; its additional factor
$X_{\beta,\gamma}(t)$ is independent of $m_1$ and is
$T^{O(1/\log T)}$ on the support of $w$.

For either AFE sign let
\begin{equation}\label{eq:completed-actual-weight}
 \mathfrak F_{d,a,\rho,q,k}^{\pm}(m_1)
 =\frac1{TS}\int_{\mathbb R}\int_{\mathbb R}
 \mathscr W^{\pm}_{dq,k,\delta m_1,d\rho}(x,t)
 e\!\left(-\frac{ax}{qk\delta m_1}\right)\dd x\dd t.
\end{equation}
Combining \eqref{eq:direct-summand}, \eqref{eq:poisson-formula}, and
\eqref{eq:poisson-jacobian}, the contribution for fixed
$d,a,\rho,q,k,m_1$, before the reciprocal phase is summed, is exactly
\begin{equation}\label{eq:completed-exact-factorization}
 \frac{d\,TS}
 {(H_0K_0L_0M_0N_0)^{1/2}H_0K_0M_0}
 a_{dq}\overline{b_k}\,
 \mathfrak F_{d,a,\rho,q,k}^{\pm}(m_1)
 e\!\left(-\frac{a\rho\overline{k_1m_1}}q\right),
\end{equation}
up to the common sign convention for the two exponentials.  The factor $d$
is precisely the numerator of the density $d/(hkm)$ and has not been hidden
in the smooth weight.

We now verify the completion hypotheses.  Fix an auxiliary
$0<\eps_0<\eps$ and invoke the localization, frequency truncation, and
divisor estimates in this proof with $\eps_0$ in place of $\eps$; its size
will be fixed after the finite completion order is known.  The support is
contained in a
fixed multiple of an interval of length
\begin{equation}\label{eq:completed-Xm}
                  X_m=\frac{M_0}{\delta}=\frac{M_0g}{d}.
\end{equation}
If this interval contains an integer, then $X_m\gg1$; otherwise the sum is
empty.  On the $x$-support one has
$x\asymp K_0L_0M_0\asymp S$, where
$S=H_0N_0+K_0L_0M_0$ is fixed throughout the block.  Put
$D_x=x\partial_x$.  Since $|d\rho|/x\ll T^{-1+\eps_0}$, the identity
\[
 \log\frac{x}{x+d\rho}=-\int_0^{d\rho}\frac{\dd v}{x+v}
\]
gives, for every fixed $j\geq0$,
\begin{equation}\label{eq:completed-phase-x-derivatives}
 \left|D_x^j\left[t\log\frac{x}{x+d\rho}\right]\right|
 \ll_j\frac{T|d\rho|}{S}\ll T^{\eps_0}.
\end{equation}
The segment $x+v$ remains comparable with $S$, also when $d\rho<0$.
Repeated differentiation of the exponential therefore costs only
$T^{O_j(\eps_0)}$.

At fixed $x$ the $t$-oscillation is independent of $m_1$, while
\begin{equation}\label{eq:completed-variable-derivatives}
 m_1\partial_{m_1}\ell_x=-\ell_x,\qquad
 m_1\partial_{m_1}n_x=0,\qquad
 m_1\partial_{m_1}\left[t\log\frac{x}{x+d\rho}\right]=0.
\end{equation}
If the moving support is instead fixed by $x=k\delta m_1L_0y$, then
\[
 \left.m_1\partial_{m_1}\right|_y
 =\left.m_1\partial_{m_1}\right|_x+D_x,
\]
so \eqref{eq:completed-phase-x-derivatives} gives the same conclusion.
In particular, this change does not differentiate the fixed block-scale
$S$.  For
\[
 \xi=\frac{x}{kmY},\qquad \chi=\frac{m(x+d\rho)}h,
\]
one has
\begin{equation}\label{eq:completed-kernel-derivatives}
 m_1\partial_{m_1}\xi=-\xi,\qquad
 m_1\partial_{m_1}\chi=\chi.
\end{equation}
Thus \eqref{eq:double-weight-bound} controls every derivative falling on
$\V^\pm$.  Derivatives of the dyadic cutoffs cost $X_m^{-1}$, and all
algebraic powers have bounded logarithmic derivatives on their supports.
Finally,
\begin{equation}\label{eq:completed-inner-phase-bound}
 \left|\frac{ax}{qkm}\right|
 \ll \frac{|a|dS}{H_0K_0M_0}
 \ll T^{\eps_0}
\end{equation}
by \eqref{eq:a-range}; hence every fixed number of derivatives of the inner
Poisson Fourier phase is also $T^{O_j(\eps_0)}$.

The $x$-support has length $O(S)$ and the $t$-support length $O(T)$.
Differentiating under the integrals in
\eqref{eq:completed-actual-weight} therefore yields, for every fixed
$j\geq0$,
\begin{equation}\label{eq:completed-weight-derivatives}
 \left|(m_1\partial_{m_1})^j
 \mathfrak F_{d,a,\rho,q,k}^{\pm}(m_1)\right|
 \ll_jT^{O_j(\eps_0)},
 \qquad
 \left|\frac{d^j}{dm_1^j}\mathfrak F_{d,a,\rho,q,k}^{\pm}(m_1)\right|
 \ll_jT^\eps X_m^{-j}.
\end{equation}
Here the localization loss $\eps_0$ is chosen sufficiently small in terms
of the final $\eps$ and the fixed number of derivatives used in completion.
No derivative order depends on $T$.

Since $(q,k_1)=1$, the coefficient in the reciprocal phase may be taken as
$c=\pm a\rho\overline{k_1}\pmod q$, for which
$(c,q)=(a\rho,q)$.  Applying \Cref{lem:reciprocal-sum} to the actual weight
gives
\begin{align}
 &\sum_{\substack{m_1\asymp M_0g/d\\(m_1,q)=1}}
 \mathfrak F_{d,a,\rho,q,k}^{\pm}(m_1)
 e\!\left(\frac{\pm a\rho\overline{k_1m_1}}q\right)\notag\\
 &\qquad\ll T^\eps
 \left\{
 \frac{M_0g}{d}q^{-1/2}(a\rho,q)^{1/2}
 +q^{1/2}(a\rho,q)^{1/2}
 \right\}.
 \label{eq:m-reciprocal}
\end{align}
Put $Q_d=H_0/d$.  By \eqref{eq:gcd-average}, restriction to
$(q,k_1)=1$ can only decrease the majorant.  Applying
\eqref{eq:gcd-average} separately to the two terms in
\eqref{eq:m-reciprocal} gives
\[
 \frac{M_0g}{d}Q_d^{-1/2}
 \sum_{q\asymp Q_d}(a\rho,q)^{1/2}
 \ll_\eps \frac{M_0g}{d}Q_d^{1/2}(a\rho)^\eps
\]
and
\[
 Q_d^{1/2}\sum_{q\asymp Q_d}(a\rho,q)^{1/2}
 \ll_\eps Q_d^{3/2}(a\rho)^\eps.
\]
\begin{samepage}
Consequently,
\begin{align}
 &\sum_{q\asymp Q_d}
 \left(\frac{M_0g}{d}q^{-1/2}+q^{1/2}\right)
 (a\rho,q)^{1/2}\notag\\
 &\qquad\ll T^\eps
 \left(\frac{M_0gH_0^{1/2}}{d^{3/2}}
       +\frac{H_0^{3/2}}{d^{3/2}}\right).
 \label{eq:completed-q-sum}
\end{align}
\end{samepage}
Indeed these are respectively $X_mQ_d^{1/2}$ and $Q_d^{3/2}$.
The factor $(a\rho)^\eps$ is absorbable because all retained dyadic
parameters are bounded by a fixed power of $T$; the auxiliary exponent is
chosen before relabelling the total loss as $T^\eps$.

Take the supremum norm of $a_{dq}$ and sum over $k$.  The sets
$\{k:(d,k)=g\}$ partition the $k$-block.  The second term of
\eqref{eq:completed-q-sum} is independent of $g$, and in the first we use
only $g\leq d$.  Therefore
\begin{align}
 &\sum_{\substack{h\asymp H_0,\,k\asymp K_0,\,m\asymp M_0\\(h,km)=d}}
 a_h\overline{b_k}\,
 \mathfrak F_{d,a,\rho,h/d,k}^{\pm}(m/\delta)
 e\!\left(\frac{\pm a\rho\overline{km/d}}{h/d}\right)
 \notag\\
 &\hspace{22mm}\ll
 T^\eps\mathcal A_{H_0}\mathcal B_{K_0}
 \left(\frac{H_0^{3/2}}{d^{3/2}}
       +\frac{M_0H_0^{1/2}}{d^{1/2}}\right).
 \label{eq:hkm-completion}
\end{align}
We finish by multiplying every scalar factor.  After the density numerator
$d$ is removed from \eqref{eq:completed-exact-factorization}, the original
Dirichlet normalization, the remaining progression density and Jacobian,
and the actual $t,x$-integral size give
\begin{equation}\label{eq:block-normalization}
 \mathcal C_0=
 \frac{TS}{(H_0K_0L_0M_0N_0)^{1/2}H_0K_0M_0}.
\end{equation}
More explicitly, the original normalization is
$(H_0K_0L_0M_0N_0)^{-1/2}$, the progression density is $d/H_0$, the
$x=kml$ Jacobian costs $(K_0M_0)^{-1}$, and the Poisson integral has size
$TS$.  For fixed $d$, the density numerator, the number of frequencies,
and the number of nonzero shifts contribute
\begin{equation}\label{eq:d-frequency-count}
 d\cdot\frac{A_0}{d}\cdot\frac{R_0}{d}
 =\frac{A_0R_0}{d},\qquad
 A_0=\frac{H_0K_0M_0}{S}T^{\eps_0},\quad
 R_0=\frac ST T^{\eps_0}.
\end{equation}
Multiplying \eqref{eq:hkm-completion}--\eqref{eq:d-frequency-count} and
extending the $d$-sum to all positive integers gives
\begin{align}
 E^{(l)}(H_0,K_0,L_0,M_0,N_0)
 \ll{}&T^\eps\mathcal C_0A_0R_0
 \mathcal A_{H_0}\mathcal B_{K_0}\notag\\
 &\times\left(
 H_0^{3/2}\sum_{d\geq1}d^{-5/2}
 +M_0H_0^{1/2}\sum_{d\geq1}d^{-3/2}\right).
 \label{eq:completed-d-series}
\end{align}
Both $d$-series converge.  Finally, by \eqref{eq:S-balance},
\begin{equation}\label{eq:completed-scalar-cancellation}
 \mathcal C_0A_0R_0
 =\frac{S}{(H_0K_0L_0M_0N_0)^{1/2}}T^{O(\eps_0)}
 \asymp T^{O(\eps_0)}.
\end{equation}
Choosing $\eps_0$ in terms of the final $\eps$ proves
\eqref{eq:block-weil}.
\end{proof}

\begin{proposition}[Nonzero-frequency estimate]
\label{prop:nonzero-global}
With the fixed $l$-Poisson orientation,
\begin{equation}\label{eq:nonzero-norm}
 \mathcal R_{\neq0}^{(l)}
 \ll_\eps T^\eps\lVert a\rVert_\infty H^{3/2}
                    \sum_{k\leq K}|b_k|.
\end{equation}
In particular, under \eqref{eq:coeff-bounds},
\begin{equation}\label{eq:nonzero-final}
                       \mathcal R_{\neq0}^{(l)}
                       \ll_\eps T^\eps KH^{3/2}.
\end{equation}
\end{proposition}

\begin{proof}
For a nonempty nonzero-frequency block, $r=d\rho\neq0$ and
\eqref{eq:r-localization} imply
\begin{equation}\label{eq:nonempty-S}
                          S\gg dT^{1-\eps}.
\end{equation}
Since $N_0\asymp S/H_0$ and $M_0N_0\ll T^{1+\eps}$,
\begin{equation}\label{eq:M-less-H}
 M_0\ll\frac{H_0T^{1+\eps}}S
      \ll\frac{H_0}{d}T^{2\eps}\leq H_0T^{2\eps}.
\end{equation}
Substitution in \eqref{eq:block-weil}, followed by the logarithmic dyadic
sums and the two AFE pieces, gives \eqref{eq:nonzero-norm}.  The divisor
bounds then give \eqref{eq:nonzero-final}.
\end{proof}

\begin{remark}\label{rem:no-minimum}
Poisson summation in another variable gives a separately valid estimate for
a different decomposition, but its zero frequency is different.  We do not
take a minimum over Poisson orientations.  The fixed orientation used in
\Cref{prop:nonzero-global} is exactly the one whose zero mode is evaluated in
the following sections.
\end{remark}
